% STATUS: conditional framework; outstanding model-specific work is explicit.
\section{P13: First Lyapunov coefficient (unfinished application)}
\label{sec:P13}
No model-specific derivative tensors or sign calculation were supplied.
\subsection{Derivative tensors}\label{slot:P13-1}
At the candidate, use the Taylor convention
$F(x_H+z)=Jz+\frac12\mathcal B(z,z)+\frac16\mathcal C(z,z,z)+O(\|z\|^4)$,
where $\mathcal B=D^2F$ and $\mathcal C=D^3F$.
Both tensors must be differentiated from the DM field, including its
rational terms and coincident offsets; DM is not a quadratic model.
\subsection{Fourier and resolvent inputs}\label{slot:P13-2}
Use P11's $Jv=\ii\omega v$, $J^*p=-\ii\omega p$, $p^*v=1$,
$\omega>0$, and its verified inverse gates. Cyclic equivariance restricts
products of modes to their sum modulo $n$; retain normalization factors
and possible aliasing when evaluating the actual tensors.
\subsection{Standard coefficient}\label{slot:P13-3}
In this convention the standard formula is \cite{KuznetsovHopf2006}
\begin{align*}
 l_1=\frac1{2\omega}\Ree\bigl\{p^*\bigl[
 &\mathcal C(v,v,\bar v)-2\mathcal B(v,J^{-1}\mathcal B(v,\bar v))\\
 &+\mathcal B(\bar v,(2\ii\omega I-J)^{-1}\mathcal B(v,v))\bigr]\bigr\}.
\end{align*}
The formula replaces a derivation of general normal-form theory; the DM
evaluation and rigorous sign or nonvanishing certificate remain required.
\subsection{Hopf conclusion}\label{slot:P13-4}
Only after P12's nonzero crossing and $l_1\ne0$ have been verified can
the nondegenerate Hopf theorem be applied. Branch stability additionally
uses P11's stable complement. No periodic branch is established here.
